% GENERATED FILE. Edit canonical lessons, then run npm run generate:books.
% canonical-source-hash: fnv1a64:6cbc18aa8125fe44
% canonical-lessons: JA-C329-joushiki, JA-C329-atarimae, JA-C329-touzen, JA-C329-junban, JA-C329-suigara

\chapter{Common Sense}
\label{ch:ja-common-sense}
\hlchaptermodality{329}

\begin{chapteropening}
\textbf{By the end of this chapter:} \emph{I can say what everyone is expected to know and do.}

Complete the last lesson of chapter 329: I can say what everyone is expected to know and do.
\end{chapteropening}

\section[\texorpdfstring{\ja{じょうしき} (jōshiki)}{じょうしき (jōshiki)}]{\ja{じょうしき} (jōshiki) — common sense}

\label{lesson:JA-C329-joushiki}



\begin{quote}
Before the new one: say the Japanese for to refrain from, then the Japanese for to cut in line.
\end{quote}

\subsection*{You'll want to know: \ja{じょうしき}}
\textbf{\ja{じょうしき}} — \emph{jōshiki} — \textquotedblleft{}common sense\textquotedblright{}.

\begin{quote}
\textbf{\ja{かれは} \ja{じょうしきが} \ja{ないと} \ja{おもいます。}} — \emph{kare wa jōshiki ga nai to omoimasu} — \textquotedblleft{}I think he has no common sense.\textquotedblright{}
\end{quote}

\begin{cousinweb}
\textbf{\ja{じょうしき}} (\emph{jōshiki}) is Sino-Japanese (kango): two Chinese characters read with their on'yomi, readings borrowed from Chinese. They are \emph{jō}, \textquotedblleft{}usual\textquotedblright{} (Middle Chinese \emph{dzyang}) and \emph{shiki}, \textquotedblleft{}know\textquotedblright{} (\emph{syik}) (KANJIDIC2; Middle Chinese in Baxter's transcription).
\end{cousinweb}

\begin{grammarlens}[title={What you've built}]
One more word for giving reasons and opinions.
\end{grammarlens}

\subsection*{Your turn}
\begin{itemize}
\raggedright
  \item \emph{Say it:} \emph{jōshiki}
  \item \emph{Say it:} \emph{jōshiki}, once more
  \item \emph{Say it:} \emph{kare wa jōshiki ga nai to omoimasu} — I think he has no common sense.
  \item \emph{Recall:} say the Japanese for to refrain from, then the Japanese for to cut in line, then say \emph{jōshiki} again
\end{itemize}

\begin{tcolorbox}[breakable,colback=teal!4,colframe=teal!35!black,title={Before you move on}]
What does \ja{じょうしき} mean? (\textquotedblleft{}Common sense\textquotedblright{}.) Say it once more.
\end{tcolorbox}
\section[\texorpdfstring{\ja{あたりまえ} (atarimae)}{あたりまえ (atarimae)}]{\ja{あたりまえ} (atarimae) — only natural; to be expected}

\label{lesson:JA-C329-atarimae}



\begin{quote}
Before the new one: say the Japanese for to cut in line, then the Japanese for common sense.
\end{quote}

\subsection*{You'll want to know: \ja{あたりまえ}}
\textbf{\ja{あたりまえ}} — \emph{atarimae} — \textquotedblleft{}only natural; to be expected\textquotedblright{}.

\begin{quote}
\textbf{\ja{あたりまえの} \ja{ことを} \ja{しました。}} — \emph{atarimae no koto o shimashita} — \textquotedblleft{}I did what anyone would have done.\textquotedblright{}
\end{quote}

A na-adjective: \ja{な} before a noun.

\begin{cousinweb}
\textbf{\ja{あたりまえ}} (\emph{atarimae}) is a native Japanese word (wago), inherited rather than borrowed from Chinese: it is built from \emph{atari}, \textquotedblleft{}hitting (ataru)\textquotedblright{} and \emph{mae}, \textquotedblleft{}front\textquotedblright{}, and its kanji were attached later as kun'yomi (JMdict).
\end{cousinweb}

\begin{grammarlens}[title={What you've built}]
One more word for giving reasons and opinions.
\end{grammarlens}

\subsection*{Your turn}
\begin{itemize}
\raggedright
  \item \emph{Say it:} \emph{atarimae}
  \item \emph{Say it:} \emph{atarimae}, once more
  \item \emph{Say it:} \emph{atarimae no koto o shimashita} — I did what anyone would have done.
  \item \emph{Recall:} say the Japanese for to cut in line, then the Japanese for common sense, then say \emph{atarimae} again
\end{itemize}

\begin{tcolorbox}[breakable,colback=teal!4,colframe=teal!35!black,title={Before you move on}]
What does \ja{あたりまえ} mean? (\textquotedblleft{}Only natural; to be expected\textquotedblright{}.) Say it once more.
\end{tcolorbox}
\section[\texorpdfstring{\ja{とうぜん} (tōzen)}{とうぜん (tōzen)}]{\ja{とうぜん} (tōzen) — naturally; as a matter of course}

\label{lesson:JA-C329-touzen}



\begin{quote}
Before the new one: say the Japanese for common sense, then the Japanese for only natural.
\end{quote}

\subsection*{You'll want to know: \ja{とうぜん}}
\textbf{\ja{とうぜん}} — \emph{tōzen} — \textquotedblleft{}naturally; as a matter of course\textquotedblright{}.

\begin{quote}
\textbf{\ja{まいにち} \ja{れんしゅうしたので、とうぜん} \ja{かちました。}} — \emph{mainichi renshū shita node, tōzen kachimashita} — \textquotedblleft{}We practised every day, so naturally we won.\textquotedblright{}
\end{quote}

\begin{cousinweb}
\textbf{\ja{とうぜん}} (\emph{tōzen}) is Sino-Japanese (kango): two Chinese characters read with their on'yomi, readings borrowed from Chinese. They are \emph{tō}, \textquotedblleft{}right\textquotedblright{} (Middle Chinese \emph{tang}) and \emph{zen}, \textquotedblleft{}so\textquotedblright{} (\emph{nyen}) (KANJIDIC2; Middle Chinese in Baxter's transcription).
\end{cousinweb}

\begin{grammarlens}[title={What you've built}]
One more word for giving reasons and opinions.
\end{grammarlens}

\subsection*{Your turn}
\begin{itemize}
\raggedright
  \item \emph{Say it:} \emph{tōzen}
  \item \emph{Say it:} \emph{tōzen}, once more
  \item \emph{Say it:} \emph{mainichi renshū shita node, tōzen kachimashita} — We practised every day, so naturally we won.
  \item \emph{Recall:} say the Japanese for common sense, then the Japanese for only natural, then say \emph{tōzen} again
\end{itemize}

\begin{tcolorbox}[breakable,colback=teal!4,colframe=teal!35!black,title={Before you move on}]
What does \ja{とうぜん} mean? (\textquotedblleft{}Naturally; as a matter of course\textquotedblright{}.) Say it once more.
\end{tcolorbox}
\section[\texorpdfstring{\ja{じゅんばん} (junban)}{じゅんばん (junban)}]{\ja{じゅんばん} (junban) — one's turn; order}

\label{lesson:JA-C329-junban}



\begin{quote}
Before the new one: say the Japanese for only natural, then the Japanese for naturally.
\end{quote}

\subsection*{You'll want to know: \ja{じゅんばん}}
\textbf{\ja{じゅんばん}} — \emph{junban} — \textquotedblleft{}one's turn; order\textquotedblright{}.

\begin{quote}
\textbf{\ja{じゅんばんが} \ja{きたら、はいって} \ja{ください。}} — \emph{junban ga kitara, haitte kudasai} — \textquotedblleft{}When your turn comes, please go in.\textquotedblright{}
\end{quote}

\begin{cousinweb}
\textbf{\ja{じゅんばん}} (\emph{junban}) is Sino-Japanese (kango): two Chinese characters read with their on'yomi, readings borrowed from Chinese. They are \emph{jun}, \textquotedblleft{}order\textquotedblright{} (Middle Chinese \emph{zywinH}) and \emph{ban}, \textquotedblleft{}turn\textquotedblright{} (\emph{phjon}) (KANJIDIC2; Middle Chinese in Baxter's transcription).
\end{cousinweb}

\begin{grammarlens}[title={What you've built}]
One more word for giving reasons and opinions.
\end{grammarlens}

\subsection*{Your turn}
\begin{itemize}
\raggedright
  \item \emph{Say it:} \emph{junban}
  \item \emph{Say it:} \emph{junban}, once more
  \item \emph{Say it:} \emph{junban ga kitara, haitte kudasai} — When your turn comes, please go in.
  \item \emph{Recall:} say the Japanese for only natural, then the Japanese for naturally, then say \emph{junban} again
\end{itemize}

\begin{tcolorbox}[breakable,colback=teal!4,colframe=teal!35!black,title={Before you move on}]
What does \ja{じゅんばん} mean? (\textquotedblleft{}One's turn; order\textquotedblright{}.) Say it once more.
\end{tcolorbox}
\section[\texorpdfstring{\ja{すいがら} (suigara)}{すいがら (suigara)}]{\ja{すいがら} (suigara) — a cigarette end}

\label{lesson:JA-C329-suigara}



\begin{quote}
Before the new one: say the Japanese for naturally, then the Japanese for one's turn.
\end{quote}

\subsection*{You'll want to know: \ja{すいがら}}
\textbf{\ja{すいがら}} — \emph{suigara} — \textquotedblleft{}a cigarette end\textquotedblright{}.

\begin{quote}
\textbf{\ja{きたないから、すいがらを} \ja{すてないで} \ja{ください。}} — \emph{kitanai kara, suigara o sutenaide kudasai} — \textquotedblleft{}It's dirty, so please do not throw away cigarette ends.\textquotedblright{}
\end{quote}

\begin{cousinweb}
\textbf{\ja{すいがら}} (\emph{suigara}) is a native Japanese word (wago), inherited rather than borrowed from Chinese: it is built from \emph{sui}, \textquotedblleft{}inhaling (suu)\textquotedblright{} and \emph{gara}, \textquotedblleft{}husk (kara)\textquotedblright{}, and its kanji were attached later as kun'yomi (JMdict).
\end{cousinweb}

\begin{grammarlens}[title={What you've built}]
One more word for giving reasons and opinions.
\end{grammarlens}

\subsection*{Your turn}
\begin{itemize}
\raggedright
  \item \emph{Say it:} \emph{suigara}
  \item \emph{Say it:} \emph{suigara}, once more
  \item \emph{Say it:} \emph{kitanai kara, suigara o sutenaide kudasai} — It's dirty, so please do not throw away cigarette ends.
  \item \emph{Recall:} say the Japanese for naturally, then the Japanese for one's turn, then say \emph{suigara} again
\end{itemize}

\begin{tcolorbox}[breakable,colback=teal!4,colframe=teal!35!black,title={Before you move on}]
What does \ja{すいがら} mean? (\textquotedblleft{}A cigarette end\textquotedblright{}.) Say it once more.
\end{tcolorbox}
